\subsection{不可分解模。初等因子}

定义3.——环A上的左模M被称为可分解的，如果它是某个由真非零子模组成的族的直和。否则，它被称为不可分解的。

因此，零模是可分解的，因为它是空子模族的直和。

设 \(\mathfrak{a}\) 是环 \(\mathbf{A}\) 的一个左理想；\(\mathbf{A} / \mathfrak{a}\) 的子模正是诸商 \(\mathfrak{b} / \mathfrak{a}\)，其中 \(\mathfrak{b}\) 是 \(\mathbf{A}\) 的包含 \(\mathfrak{a}\) 的理想（I，第 41 页，定理 4）；若 \(\mathfrak{b}\) 与 \(\mathfrak{c}\) 是 \(\mathbf{A}\) 的包含 \(\mathfrak{a}\) 的两个理想，则模 \(\mathbf{A} / \mathfrak{a}\) 是子模 \(\mathfrak{b} / \mathfrak{a}\) 与 \(\mathfrak{c} / \mathfrak{a}\) 的直和，当且仅当 \(\mathbf{A} = \mathfrak{b} + \mathfrak{c}\) 且 \(\mathfrak{b} \cap \mathfrak{c} = \mathfrak{a}\)。因此：

引理2. —— 模 A/a 不可分解当且仅当 a ≠ A，并且不存在 A 的理想对 (b, c)，它们均不同于 A 和 a，使得 A = b + c 且 b ∩ c = a。

命题7.——设 A 是一个交换环，设 p 是 A 的一个素理想（I，第 117 页，定义 3），并设 q 是 A 的包含在 p 中的一个理想。假设对每个 \(x \in p\) 存在一个整数 n \textgreater{} 0 使得 \(x^{n} \in q\)。则 A-模 A/q 是不可分解的。

设 \(\mathfrak{b}\) 与 \(\mathfrak{c}\) 是 \(\mathbf{A}\) 的两个理想，使得 \(\mathbf{A} = \mathfrak{b} + \mathfrak{c}\) 且 \(\mathfrak{b} \cap \mathfrak{c} = \mathfrak{q}\)。则 \(\mathfrak{bc} \subset \mathfrak{b} \cap \mathfrak{c} = \mathfrak{q} \subset \mathfrak{p}\)；若 \(x \notin \mathfrak{p}\) 且 \(x \in \mathfrak{c}\)，则 \(xb \subset \mathfrak{p}\)，于是 \(\mathfrak{b} \subset \mathfrak{p}\)（I，第 116 页，命题 4）；因而或者 \(\mathfrak{b} \subset \mathfrak{p}\)，或者 \(\mathfrak{c} \subset \mathfrak{p}\)。例如设 \(\mathfrak{c} \subset \mathfrak{p}\)，于是 \(\mathfrak{b} + \mathfrak{p} = \mathbf{A}\)；则存在 \(x \in \mathfrak{b}\) 与 \(y \in \mathfrak{p}\) 使得 \(1 = x + y\)；设 \(n \in \mathbf{N}\) 使得 \(y^n \in \mathfrak{q}\)；则 \(1 = (x + y)^n\)，故 \(1 \in x\mathbf{A} + y^n\mathbf{A} \subset \mathfrak{b} + \mathfrak{q} \subset \mathfrak{b}\)，于是 \(\mathfrak{b} = \mathbf{A}\)。引理 2 现在表明 \(\mathbf{A}/\mathfrak{q}\) 是不可分解的。

现在假设 A 是一个主理想整环；由 VII 第 2 页命题 2，A 的素理想就是诸理想 \((p)\)（其中 p 是 A 的不可约元）以及理想 0；由前一命题，模 A 与 \(\mathbf{A}/(p^{n})\)（其中 p 不可约且 n \textgreater{} 0）都是不可分解的。由于每个循环模都是此类模的直和（VII，第3页，命题4），且每个有限生成的A-模都是循环模的直和（VII，第19页，定理2），我们推得：

命题8.——设A为主理想整环，且M为有限生成的A-模。

\begin{enumerate}
\def\labelenumi{\alph{enumi})}
\item
  M 是不可分解的，当且仅当它同构于 A，或同构于形如 \(\mathbf{A} / (p^n)\) 的模，其中 \(p\) 是 A 的一个不可约元，且 \(n > 0\) 是整数。
\item
  M 是有限个不可分解子模族的直和。
\end{enumerate}

上述命题的(b)部分可以更精确地表述如下：

命题9. --- 设 A 为主理想整环，设 P 为 A 的不可约元的一个代表系，并设 M 为有限生成的 A-模。则存在由 M 唯一确定的正整数 \(m(0)\) 与 \(m(p^{n})\)（\(p \in P\)，n \textgreater{} 0），且除有限多个外均为零，使得 M 同构于 \(\mathbf{A}^{m(0)}\) 与诸 \((\mathbf{A}/(p^{n}))^{m(p^{n})}\)（\(p \in P\)，n \textgreater{} 0）的直和。

整数 \(m(0)\) 与 \(m(p^n)\)（\(p \in P\)，\(n > 0\)）的存在性由命题 8 得出。整数 \(m(0)\) 是唯一确定的：它是 \(M\) 对其挠子模的商自由模的秩。最后，对每个 \(p\)，\(M\) 的相应准素分量同构于诸 \((A/(p^n))^{m(p^n)}\) 的直和；由于理想族 \((p^n)\)（\(n \geqslant 1\)）按包含关系全序，故 \(m(p^n)\) 的唯一性由 VII 第 16 页命题 2 的推论得出。

定义4. --- 在命题 9 的记号下，那些理想 \((p^{n})\)（\(p \in P\)，\(n \geqslant 1\) 为整数），若满足 \(m(p^{n}) > 0\)，则称为模 M 的初等因子，而诸整数 \(m(p^{n})\) 称为它们的重数；若整数 \(m(0)\) \textgreater{} 0，则称之为初等因子 0 的重数。

正如对不变因子那样（VII，第 19 页，定义 1），当 A = Z 或 A = K{[}X{]}（K 为交换域）时，理想 \((p^{n})\) 的典范生成元（正整数或首一多项式）（滥用术语）也称为有限生成模 M 的初等因子。

注记。--- 1) 若 M 是有限阿贝尔群，则其结构可通过写出其初等因子（每个按其重数重复）来描述。例如，当 M 同构于两个群 \(\mathbf{Z}/(2)\)、一个群 \(\mathbf{Z}/(2^{2})\)、两个群 \(\mathbf{Z}/(3^{3})\) 和一个群 \(\mathbf{Z}/(5^{2})\) 的乘积时，我们就说 M “属于类型 (2, 2, 4, 27, 27, 25)”（或说它是“一个 (2, 2, 4, 27, 27, 25) 型群”）.

\begin{enumerate}
\def\labelenumi{\arabic{enumi})}
\setcounter{enumi}{1}
\tightlist
\item
  如果一个有限生成挠模 \(\mathbf{M}\)（定义在主理想整环 \(\mathbf{A}\) 上）给定为同构于 \(\mathbf{A} / (a_i)\) 的循环模的直和（特别地，若已知 \(\mathbf{M}\) 的不变因子），则 \(\mathbf{M}\) 的初等因子及其重数可以这样确定：注意到 \(\mathbf{A} / (a)\) 同构于诸 \(\mathbf{A} / (p^{n(p)})\) 的乘积，其中 \(a = \varepsilon \prod_{p \in \mathbb{P}} p^{n(p)}\) 是 \(a\) 的不可约因子分解
\end{enumerate}

因子（VII，第 3 页）。让我们以乘法群 G(464 600) 为例进行研究，其中 G(n) 表示乘法群 \((\mathbf{Z}/n\mathbf{Z})^{*}\)（VII，第 12 页）。由于 \(464\;600 = 2^{3}\cdot5^{2}\cdot23\cdot101\)，该群同构于群 G( \(2^{3}\) )、G( \(5^{2}\) )、G(23) 与 G(101) 的乘积（VII，第 13 页，定理 3）；而后三个群分别是阶为 20、22 和 100 的循环群，G( \(2^{3}\) ) 是两个 2 阶循环群的乘积（同上）；由于 \(20 = 2^{2}\cdot5\)、22 = \(2\cdot 11\) 以及 \(100 = 2^{2}\cdot5^{2}\)，群 G(464 600) 属于类型 (2, 2, 2, \(2^{2}\), \(2^{2}\), 5, \(5^{2}\), 11).

\begin{enumerate}
\def\labelenumi{\arabic{enumi})}
\setcounter{enumi}{2}
\tightlist
\item
  为了计算初等因子已知的挠模的不变因子，我们仍然利用如下事实：若 \(a_{i}\) 是 A 中两两互素的元素，则乘积 \(\prod_{i}\mathbf{A}/(a_{i})\) 是一个循环模，同构于
\end{enumerate}

\(A/(a_{1}a_{2}\ldots a_{n})\)（VII，第 3 页，命题 4）。让我们通过考察群 G(464 600) = M 的例子来说明这一方法：写出 M 的初等因子 \(p^{n}\)，把同一不可约元 p 的幂写在同一行，从指数最大的开始；必要时补入 1，把这些行延长成等长的行：

\[
\begin{array}{c c c c c} 2 ^ {2}, & 2 ^ {2}, & 2, & 2, & 2 \\ 5 ^ {2}, & 5, & 1, & 1, & 1 \\ 11, & 1, & 1, & 1, & 1. \end{array}
\]

则不变因子就是同一列中元素的乘积：1100、20、2、2、2。事实上，由 VII 第 3 页命题 4，M 同构于阶为 1100、20、2、2、2 的循环群的乘积；由于这些阶中每一个都是下一个的倍数，因此这些就是 M 的不变因子（VII，第 22 页，定理 3）。

\begin{enumerate}
\def\labelenumi{\arabic{enumi})}
\setcounter{enumi}{3}
\tightlist
\item
  一个 A-模称为单模（I，第 37 页），如果它非零且除自身和 0 外没有其他子模；它此时必为循环模，因此是有限生成的，并且是不可分解的；由于模 \(\mathbf{A} / (p^n)\) 对 \(n \neq 1\) 不是单模，而模 \(\mathbf{A} / (p)\) 是单模，且 \(\mathbf{A}\) 只有当环 \(\mathbf{A}\) 是域时才是单的，我们推断出单 A-模为：
\end{enumerate}

\begin{enumerate}
\def\labelenumi{\alph{enumi})}
\item
  秩为 1 的自由模，当 \(\mathbf{A}\) 是域时；
\item
  同构于商 \(\mathbf{A} / (p)\) 的模，其中 \(p\) 是 \(\mathbf{A}\) 的不可约元，当 \(\mathbf{A}\) 不是域时。
\end{enumerate}

